\chapter{Unit 2: Conditionals, Iterations, Functions and Recursion}

\section{Section A: Multiple Choice Practice Questions (60 Curated Problems)}

\noindent This section contains 60 rigorous multiple-choice questions extracted from university course examination assessments. Review each question, test edge-case syntax, and verify code execution mechanics.
\vspace{3mm}

\noindent\textbf{Q1.} \hfill {\small\color{LPUOrange}\textbf{[if statement | Easy]}}\\[2pt]
Which keyword is used to begin a conditional statement in Python?

\begin{enumerate}[label=(\alph*),topsep=2pt,itemsep=1pt,partopsep=0pt,parsep=0pt]
    \item \relax \texttt{when}
    \item \relax \texttt{if}
    \item \relax \texttt{check}
    \item \relax \texttt{cond}
\end{enumerate}
\vspace{2mm}

\noindent\textbf{Q2.} \hfill {\small\color{LPUOrange}\textbf{[two way if - else | Easy]}}\\[2pt]
In a two-way \texttt{if-else} statement, when does the \texttt{else} block execute?

\begin{enumerate}[label=(\alph*),topsep=2pt,itemsep=1pt,partopsep=0pt,parsep=0pt]
    \item \relax When the \texttt{if} condition is false
    \item \relax When the \texttt{if} condition is true
    \item \relax Only when the condition evaluates to a number greater than zero, considering all edge cases carefully
    \item \relax Always, regardless of the condition
\end{enumerate}
\vspace{2mm}

\noindent\textbf{Q3.} \hfill {\small\color{LPUOrange}\textbf{[nested if and multi-way if-elif-else statement | Easy]}}\\[2pt]
Which keyword is used for a multi-way branch to test additional conditions in Python?

\begin{enumerate}[label=(\alph*),topsep=2pt,itemsep=1pt,partopsep=0pt,parsep=0pt]
    \item \relax \texttt{elseif}
    \item \relax \texttt{elsif}
    \item \relax \texttt{elif}
    \item \relax \texttt{orif}
\end{enumerate}
\vspace{2mm}

\noindent\textbf{Q4.} \hfill {\small\color{LPUOrange}\textbf{[for loop | Easy]}}\\[2pt]
What does the \texttt{range(5)} function generate when used in a \texttt{for} loop?

\begin{enumerate}[label=(\alph*),topsep=2pt,itemsep=1pt,partopsep=0pt,parsep=0pt]
    \item \relax Numbers \texttt{0, 1, 2, 3, 4}
    \item \relax Numbers \texttt{0, 1, 2, 3, 4, 5}
    \item \relax Numbers \texttt{1, 2, 3, 4, 5}
    \item \relax Numbers \texttt{1, 2, 3, 4}
\end{enumerate}
\vspace{2mm}

\noindent\textbf{Q5.} \hfill {\small\color{LPUOrange}\textbf{[while loop | Easy]}}\\[2pt]
A \texttt{while} loop in Python continues to execute as long as its condition is:

\begin{enumerate}[label=(\alph*),topsep=2pt,itemsep=1pt,partopsep=0pt,parsep=0pt]
    \item \relax \texttt{True}
    \item \relax \texttt{False}
    \item \relax \texttt{zero}
    \item \relax \texttt{None}
\end{enumerate}
\vspace{2mm}

\noindent\textbf{Q6.} \hfill {\small\color{LPUOrange}\textbf{[break and continue | Easy]}}\\[2pt]
What does the \texttt{break} statement do inside a loop?

\begin{enumerate}[label=(\alph*),topsep=2pt,itemsep=1pt,partopsep=0pt,parsep=0pt]
    \item \relax Pauses the loop temporarily until resumed by the interpreter automatically
    \item \relax Restarts the loop from the beginning
    \item \relax Skips to the next iteration of the loop
    \item \relax Exits the loop immediately
\end{enumerate}
\vspace{2mm}

\noindent\textbf{Q7.} \hfill {\small\color{LPUOrange}\textbf{[break and continue | Easy]}}\\[2pt]
What does the \texttt{continue} statement do inside a loop?

\begin{enumerate}[label=(\alph*),topsep=2pt,itemsep=1pt,partopsep=0pt,parsep=0pt]
    \item \relax Stops the program execution
    \item \relax Exits the loop completely
    \item \relax Repeats the current iteration
    \item \relax Skips the rest of the current iteration and moves to the next
\end{enumerate}
\vspace{2mm}

\noindent\textbf{Q8.} \hfill {\small\color{LPUOrange}\textbf{[nested loops | Easy]}}\\[2pt]
What is a nested loop?

\begin{enumerate}[label=(\alph*),topsep=2pt,itemsep=1pt,partopsep=0pt,parsep=0pt]
    \item \relax A loop that runs backwards
    \item \relax A loop that never ends
    \item \relax A loop with only one iteration
    \item \relax A loop placed inside another loop
\end{enumerate}
\vspace{2mm}

\noindent\textbf{Q9.} \hfill {\small\color{LPUOrange}\textbf{[random numbers | Easy]}}\\[2pt]
Which module in Python is commonly used to generate random numbers?

\begin{enumerate}[label=(\alph*),topsep=2pt,itemsep=1pt,partopsep=0pt,parsep=0pt]
    \item \relax \texttt{numbers}
    \item \relax \texttt{random}
    \item \relax \texttt{rand}
    \item \relax \texttt{math}
\end{enumerate}
\vspace{2mm}

\noindent\textbf{Q10.} \hfill {\small\color{LPUOrange}\textbf{[random numbers | Easy]}}\\[2pt]
Which function returns a random integer between two given values (inclusive)?

\begin{enumerate}[label=(\alph*),topsep=2pt,itemsep=1pt,partopsep=0pt,parsep=0pt]
    \item \relax \texttt{random.randint(a, b)}
    \item \relax \texttt{random.range()}
    \item \relax \texttt{random.choice()}
    \item \relax \texttt{random.random()}
\end{enumerate}
\vspace{2mm}

\noindent\textbf{Q11.} \hfill {\small\color{LPUOrange}\textbf{[function calls | Easy]}}\\[2pt]
How do you call a function named \texttt{greet} in Python?

\begin{enumerate}[label=(\alph*),topsep=2pt,itemsep=1pt,partopsep=0pt,parsep=0pt]
    \item \relax \texttt{greet()}
    \item \relax \texttt{function greet}
    \item \relax \texttt{greet[]}
    \item \relax \texttt{call greet}
\end{enumerate}
\vspace{2mm}

\noindent\textbf{Q12.} \hfill {\small\color{LPUOrange}\textbf{[type conversion and coercion | Easy]}}\\[2pt]
Which function converts a string \texttt{"25"} into an integer?

\begin{enumerate}[label=(\alph*),topsep=2pt,itemsep=1pt,partopsep=0pt,parsep=0pt]
    \item \relax \texttt{int("25")}
    \item \relax \texttt{chr(25)}
    \item \relax \texttt{float("25")}
    \item \relax \texttt{str(25)}
\end{enumerate}
\vspace{2mm}

\noindent\textbf{Q13.} \hfill {\small\color{LPUOrange}\textbf{[type conversion and coercion | Easy]}}\\[2pt]
What is the result of \texttt{float(7)} in Python?

\begin{enumerate}[label=(\alph*),topsep=2pt,itemsep=1pt,partopsep=0pt,parsep=0pt]
    \item \relax \texttt{70}
    \item \relax \texttt{7.0}
    \item \relax \texttt{"7"}
    \item \relax \texttt{7}
\end{enumerate}
\vspace{2mm}

\noindent\textbf{Q14.} \hfill {\small\color{LPUOrange}\textbf{[math functions | Easy]}}\\[2pt]
Which function from the \texttt{math} module returns the square root of a number?

\begin{enumerate}[label=(\alph*),topsep=2pt,itemsep=1pt,partopsep=0pt,parsep=0pt]
    \item \relax \texttt{math.sqrt()}
    \item \relax \texttt{math.pow()}
    \item \relax \texttt{math.root()}
    \item \relax \texttt{math.square()}
\end{enumerate}
\vspace{2mm}

\noindent\textbf{Q15.} \hfill {\small\color{LPUOrange}\textbf{[math functions | Easy]}}\\[2pt]
What does \texttt{math.ceil(4.2)} return?

\begin{enumerate}[label=(\alph*),topsep=2pt,itemsep=1pt,partopsep=0pt,parsep=0pt]
    \item \relax \texttt{4}
    \item \relax \texttt{4.0}
    \item \relax \texttt{4.2}
    \item \relax \texttt{5}
\end{enumerate}
\vspace{2mm}

\noindent\textbf{Q16.} \hfill {\small\color{LPUOrange}\textbf{[adding new function | Easy]}}\\[2pt]
Which keyword is used to define a new function in Python?

\begin{enumerate}[label=(\alph*),topsep=2pt,itemsep=1pt,partopsep=0pt,parsep=0pt]
    \item \relax \texttt{def}
    \item \relax \texttt{func}
    \item \relax \texttt{define}
    \item \relax \texttt{function}
\end{enumerate}
\vspace{2mm}

\noindent\textbf{Q17.} \hfill {\small\color{LPUOrange}\textbf{[parameters and argument | Easy]}}\\[2pt]
In the function definition \texttt{def add(x, y):}, what are \texttt{x} and \texttt{y} called?

\begin{enumerate}[label=(\alph*),topsep=2pt,itemsep=1pt,partopsep=0pt,parsep=0pt]
    \item \relax Parameters
    \item \relax Keywords
    \item \relax Variables
    \item \relax Arguments
\end{enumerate}
\vspace{2mm}

\noindent\textbf{Q18.} \hfill {\small\color{LPUOrange}\textbf{[recursion and its use | Easy]}}\\[2pt]
What is recursion in programming?

\begin{enumerate}[label=(\alph*),topsep=2pt,itemsep=1pt,partopsep=0pt,parsep=0pt]
    \item \relax A function that calls itself
    \item \relax A loop that runs forever
    \item \relax A function with no parameters
    \item \relax A function that returns nothing
\end{enumerate}
\vspace{2mm}

\noindent\textbf{Q19.} \hfill {\small\color{LPUOrange}\textbf{[recursion and its use | Easy]}}\\[2pt]
What is required in every recursive function to prevent infinite recursion?

\begin{enumerate}[label=(\alph*),topsep=2pt,itemsep=1pt,partopsep=0pt,parsep=0pt]
    \item \relax A print statement that displays the current value at each level of the recursive call chain
    \item \relax A base case
    \item \relax A global variable
    \item \relax An infinite loop
\end{enumerate}
\vspace{2mm}

\noindent\textbf{Q20.} \hfill {\small\color{LPUOrange}\textbf{[conditional expressions | Easy]}}\\[2pt]
What is the value of the conditional expression \texttt{"Yes" if 5 > 3 else "No"}?

\begin{enumerate}[label=(\alph*),topsep=2pt,itemsep=1pt,partopsep=0pt,parsep=0pt]
    \item \relax \texttt{"No"}
    \item \relax \texttt{"Yes"}
    \item \relax \texttt{True}
    \item \relax \texttt{False}
\end{enumerate}
\vspace{2mm}

\noindent\textbf{Q21.} \hfill {\small\color{LPUOrange}\textbf{[conditional expressions | Medium]}}\\[2pt]
What is the value assigned to \texttt{result} after executing the following code? 

\begin{pycode}
x = 7\\
result = "even" if x \% 2 == 0 else "odd"
\end{pycode}

\begin{enumerate}[label=(\alph*),topsep=2pt,itemsep=1pt,partopsep=0pt,parsep=0pt]
    \item \relax \texttt{7}
    \item \relax \texttt{True}
    \item \relax \texttt{"even"}
    \item \relax \texttt{"odd"}
\end{enumerate}
\vspace{2mm}

\noindent\textbf{Q22.} \hfill {\small\color{LPUOrange}\textbf{[if statement | Medium]}}\\[2pt]
How many times will \texttt{"Hi"} be printed? 

\begin{pycode}
n = 10\\
if n > 5:\\
if n < 20:\\
print("Hi")
\end{pycode}

\begin{enumerate}[label=(\alph*),topsep=2pt,itemsep=1pt,partopsep=0pt,parsep=0pt]
    \item \relax \texttt{0}
    \item \relax \texttt{1}
    \item \relax \texttt{10}
    \item \relax \texttt{2}
\end{enumerate}
\vspace{2mm}

\noindent\textbf{Q23.} \hfill {\small\color{LPUOrange}\textbf{[two way if - else | Medium]}}\\[2pt]
What does the following code output? 

\begin{pycode}
age = 18\\
if age >= 18:\\
print("Adult")\\
else:\\
print("Minor")
\end{pycode}

\begin{enumerate}[label=(\alph*),topsep=2pt,itemsep=1pt,partopsep=0pt,parsep=0pt]
    \item \relax \texttt{Adult}
    \item \relax Nothing is printed
    \item \relax \texttt{AdultMinor}
    \item \relax \texttt{Minor}
\end{enumerate}
\vspace{2mm}

\noindent\textbf{Q24.} \hfill {\small\color{LPUOrange}\textbf{[nested if and multi-way if-elif-else statement | Medium]}}\\[2pt]
What is printed by this code? 

\begin{pycode}
marks = 75\\
if marks >= 90:\\
print("A")\\
elif marks >= 70:\\
print("B")\\
elif marks >= 50:\\
print("C")\\
else:\\
print("F")
\end{pycode}

\begin{enumerate}[label=(\alph*),topsep=2pt,itemsep=1pt,partopsep=0pt,parsep=0pt]
    \item \relax \texttt{F}
    \item \relax \texttt{B}
    \item \relax \texttt{C}
    \item \relax \texttt{A}
\end{enumerate}
\vspace{2mm}

\noindent\textbf{Q25.} \hfill {\small\color{LPUOrange}\textbf{[for loop | Medium]}}\\[2pt]
What is the output of the following loop? 

\begin{pycode}
total = 0\\
for i in range(1, 5):\\
total += i\\
print(total)
\end{pycode}

\begin{enumerate}[label=(\alph*),topsep=2pt,itemsep=1pt,partopsep=0pt,parsep=0pt]
    \item \relax \texttt{5}
    \item \relax \texttt{6}
    \item \relax \texttt{10}
    \item \relax \texttt{15}
\end{enumerate}
\vspace{2mm}

\noindent\textbf{Q26.} \hfill {\small\color{LPUOrange}\textbf{[while loop | Medium]}}\\[2pt]
How many times does the \texttt{while} loop body execute? 

\begin{pycode}
x = 16\\
count = 0\\
while x > 1:\\
x = x // 2\\
count += 1\\
print(count)
\end{pycode}

\begin{enumerate}[label=(\alph*),topsep=2pt,itemsep=1pt,partopsep=0pt,parsep=0pt]
    \item \relax \texttt{16}
    \item \relax \texttt{4}
    \item \relax \texttt{5}
    \item \relax \texttt{3}
\end{enumerate}
\vspace{2mm}

\noindent\textbf{Q27.} \hfill {\small\color{LPUOrange}\textbf{[nested loops | Medium]}}\\[2pt]
What is the output of the following nested loops? 

\begin{pycode}
count = 0\\
for i in range(3):\\
for j in range(2):\\
count += 1\\
print(count)
\end{pycode}

\begin{enumerate}[label=(\alph*),topsep=2pt,itemsep=1pt,partopsep=0pt,parsep=0pt]
    \item \relax \texttt{3}
    \item \relax \texttt{5}
    \item \relax \texttt{6}
    \item \relax \texttt{2}
\end{enumerate}
\vspace{2mm}

\noindent\textbf{Q28.} \hfill {\small\color{LPUOrange}\textbf{[break and continue | Medium]}}\\[2pt]
What does this code print? 

\begin{pycode}
for i in range(1, 6):\\
if i == 3:\\
continue\\
if i == 5:\\
break\\
print(i, end=" ")
\end{pycode}

\begin{enumerate}[label=(\alph*),topsep=2pt,itemsep=1pt,partopsep=0pt,parsep=0pt]
    \item \relax \texttt{1 2 3 4 5}
    \item \relax \texttt{1 2 4 5}
    \item \relax \texttt{1 2 4}
    \item \relax \texttt{1 2 3 4}
\end{enumerate}
\vspace{2mm}

\noindent\textbf{Q29.} \hfill {\small\color{LPUOrange}\textbf{[random numbers | Medium]}}\\[2pt]
Which function call returns a random integer that could be any value in the set { 1 , 2 , 3 , 4 , 5 , 6 } ?

\begin{enumerate}[label=(\alph*),topsep=2pt,itemsep=1pt,partopsep=0pt,parsep=0pt]
    \item \relax \texttt{random.range(1, 6)}
    \item \relax \texttt{random.random(1, 6)}
    \item \relax \texttt{random.randint(1, 6)}
    \item \relax \texttt{random.randrange(1, 6)}
\end{enumerate}
\vspace{2mm}

\noindent\textbf{Q30.} \hfill {\small\color{LPUOrange}\textbf{[function calls | Medium]}}\\[2pt]
What is printed by the following code? 

\begin{pycode}
def square(n):\\
return n * n\\
~\\
print(square(square(2)))
\end{pycode}

\begin{enumerate}[label=(\alph*),topsep=2pt,itemsep=1pt,partopsep=0pt,parsep=0pt]
    \item \relax \texttt{8}
    \item \relax \texttt{2}
    \item \relax \texttt{16}
    \item \relax \texttt{4}
\end{enumerate}
\vspace{2mm}

\noindent\textbf{Q31.} \hfill {\small\color{LPUOrange}\textbf{[type conversion and coercion | Medium]}}\\[2pt]
What is the result of the following expression? 

\begin{pycode}
int("12") + int(3.9)
\end{pycode}

\begin{enumerate}[label=(\alph*),topsep=2pt,itemsep=1pt,partopsep=0pt,parsep=0pt]
    \item \relax \texttt{16}
    \item \relax \texttt{"123"}
    \item \relax \texttt{15}
    \item \relax \texttt{15.9}
\end{enumerate}
\vspace{2mm}

\noindent\textbf{Q32.} \hfill {\small\color{LPUOrange}\textbf{[type conversion and coercion | Medium]}}\\[2pt]
What is the output of the following code? 

\begin{pycode}
print(3 + True + 2.0)
\end{pycode}

\begin{enumerate}[label=(\alph*),topsep=2pt,itemsep=1pt,partopsep=0pt,parsep=0pt]
    \item \relax \texttt{5}
    \item \relax \texttt{6}
    \item \relax \texttt{6.0}
    \item \relax \texttt{TypeError}
\end{enumerate}
\vspace{2mm}

\noindent\textbf{Q33.} \hfill {\small\color{LPUOrange}\textbf{[math functions | Medium]}}\\[2pt]
What does the following code print? 

\begin{pycode}
import math\\
print(math.ceil(4.1) + math.floor(4.9))
\end{pycode}

\begin{enumerate}[label=(\alph*),topsep=2pt,itemsep=1pt,partopsep=0pt,parsep=0pt]
    \item \relax \texttt{4}
    \item \relax \texttt{9}
    \item \relax \texttt{10}
    \item \relax \texttt{8}
\end{enumerate}
\vspace{2mm}

\noindent\textbf{Q34.} \hfill {\small\color{LPUOrange}\textbf{[math functions | Medium]}}\\[2pt]
What is the value of the following expression? 

\begin{pycode}
import math\\
math.sqrt(math.pow(3, 2) + math.pow(4, 2))
\end{pycode}

\begin{enumerate}[label=(\alph*),topsep=2pt,itemsep=1pt,partopsep=0pt,parsep=0pt]
    \item \relax \texttt{25.0}
    \item \relax \texttt{5.0}
    \item \relax \texttt{12.0}
    \item \relax \texttt{7.0}
\end{enumerate}
\vspace{2mm}

\noindent\textbf{Q35.} \hfill {\small\color{LPUOrange}\textbf{[adding new function | Medium]}}\\[2pt]
What does the following program print? 

\begin{pycode}
def greet(name):\\
print("Hello", name)\\
~\\
x = greet("Sam")\\
print(x)
\end{pycode}

\begin{enumerate}[label=(\alph*),topsep=2pt,itemsep=1pt,partopsep=0pt,parsep=0pt]
    \item \relax \texttt{None} only
    \item \relax \texttt{Hello Sam} only
    \item \relax \texttt{Hello Sam} then \texttt{Sam}
    \item \relax \texttt{Hello Sam} then \texttt{None}
\end{enumerate}
\vspace{2mm}

\noindent\textbf{Q36.} \hfill {\small\color{LPUOrange}\textbf{[parameters and argument | Medium]}}\\[2pt]
What is the output of this code using default and keyword arguments? 

\begin{pycode}
def power(base, exp=2):\\
return base ** exp\\
~\\
print(power(3), power(2, exp=3))
\end{pycode}

\begin{enumerate}[label=(\alph*),topsep=2pt,itemsep=1pt,partopsep=0pt,parsep=0pt]
    \item \relax \texttt{8 9}
    \item \relax \texttt{9 6}
    \item \relax \texttt{6 8}
    \item \relax \texttt{9 8}
\end{enumerate}
\vspace{2mm}

\noindent\textbf{Q37.} \hfill {\small\color{LPUOrange}\textbf{[parameters and argument | Medium]}}\\[2pt]
What does the following code print? 

\begin{pycode}
def info(a, b, c):\\
return a - b - c\\
~\\
print(info(c=1, a=10, b=3))
\end{pycode}

\begin{enumerate}[label=(\alph*),topsep=2pt,itemsep=1pt,partopsep=0pt,parsep=0pt]
    \item \relax \texttt{8}
    \item \relax \texttt{4}
    \item \relax \texttt{6}
    \item \relax \texttt{12}
\end{enumerate}
\vspace{2mm}

\noindent\textbf{Q38.} \hfill {\small\color{LPUOrange}\textbf{[recursion and its use | Medium]}}\\[2pt]
What value does \texttt{fact(4)} return? 

\begin{pycode}
def fact(n):\\
if n <= 1:\\
return 1\\
return n * fact(n - 1)
\end{pycode}

\begin{enumerate}[label=(\alph*),topsep=2pt,itemsep=1pt,partopsep=0pt,parsep=0pt]
    \item \relax \texttt{4}
    \item \relax \texttt{24}
    \item \relax \texttt{12}
    \item \relax \texttt{10}
\end{enumerate}
\vspace{2mm}

\noindent\textbf{Q39.} \hfill {\small\color{LPUOrange}\textbf{[recursion and its use | Medium]}}\\[2pt]
What is printed by the following recursive function call \texttt{mystery(3)}? 

\begin{pycode}
def mystery(n):\\
if n == 0:\\
return\\
print(n, end=" ")\\
mystery(n - 1)
\end{pycode}

\begin{enumerate}[label=(\alph*),topsep=2pt,itemsep=1pt,partopsep=0pt,parsep=0pt]
    \item \relax \texttt{0 1 2 3}
    \item \relax \texttt{3 2 1 0}
    \item \relax \texttt{3 2 1}
    \item \relax \texttt{1 2 3}
\end{enumerate}
\vspace{2mm}

\noindent\textbf{Q40.} \hfill {\small\color{LPUOrange}\textbf{[for loop | Medium]}}\\[2pt]
What is the output of the following code that uses \texttt{range} with a step? 

\begin{pycode}
result = []\\
for i in range(10, 0, -3):\\
result.append(i)\\
print(result)
\end{pycode}

\begin{enumerate}[label=(\alph*),topsep=2pt,itemsep=1pt,partopsep=0pt,parsep=0pt]
    \item \relax \texttt{[9, 6, 3]}
    \item \relax \texttt{[10, 7, 4, 1]}
    \item \relax \texttt{[10, 7, 4, 1, -2]}
    \item \relax \texttt{[10, 7, 4]}
\end{enumerate}
\vspace{2mm}

\noindent\textbf{Q41.} \hfill {\small\color{LPUOrange}\textbf{[conditional expressions | Hard]}}\\[2pt]
What is the output of the following code? 

\begin{pycode}
x = 0\\
y = 5\\
result = "A" if x else "B" if y else "C"\\
print(result)
\end{pycode}

\begin{enumerate}[label=(\alph*),topsep=2pt,itemsep=1pt,partopsep=0pt,parsep=0pt]
    \item \relax B
    \item \relax A
    \item \relax C
    \item \relax Raises SyntaxError
\end{enumerate}
\vspace{2mm}

\noindent\textbf{Q42.} \hfill {\small\color{LPUOrange}\textbf{[if statement | Hard]}}\\[2pt]
Consider: 

\begin{pycode}
values = [0, '', None, [], 'False', 0.0]\\
count = 0\\
for v in values:\\
if v:\\
count += 1\\
print(count)
\end{pycode}

 What is printed?

\begin{enumerate}[label=(\alph*),topsep=2pt,itemsep=1pt,partopsep=0pt,parsep=0pt]
    \item \relax 1
    \item \relax 0
    \item \relax 2
    \item \relax 6
\end{enumerate}
\vspace{2mm}

\noindent\textbf{Q43.} \hfill {\small\color{LPUOrange}\textbf{[two way if - else | Hard]}}\\[2pt]
What does this function return for \texttt{f(-4)}? 

\begin{pycode}
def f(n):\\
if n \% 2 == 0:\\
return n // 2\\
else:\\
return 3 * n + 1
\end{pycode}

\begin{enumerate}[label=(\alph*),topsep=2pt,itemsep=1pt,partopsep=0pt,parsep=0pt]
    \item \relax 2
    \item \relax -2
    \item \relax -1
    \item \relax -11
\end{enumerate}
\vspace{2mm}

\noindent\textbf{Q44.} \hfill {\small\color{LPUOrange}\textbf{[nested if and multi-way if-elif-else statement | Hard]}}\\[2pt]
What is the output? 

\begin{pycode}
score = 85\\
if score >= 90:\\
grade = 'A'\\
elif score >= 80:\\
if score >= 85:\\
grade = 'B+'\\
else:\\
grade = 'B'\\
elif score >= 70:\\
grade = 'C'\\
else:\\
grade = 'F'\\
print(grade)
\end{pycode}

\begin{enumerate}[label=(\alph*),topsep=2pt,itemsep=1pt,partopsep=0pt,parsep=0pt]
    \item \relax C
    \item \relax B
    \item \relax A
    \item \relax B+
\end{enumerate}
\vspace{2mm}

\noindent\textbf{Q45.} \hfill {\small\color{LPUOrange}\textbf{[for loop | Hard]}}\\[2pt]
What is the final value of \texttt{total}? 

\begin{pycode}
total = 0\\
for i in range(2, 20, 3):\\
total += i\\
print(total)
\end{pycode}

\begin{enumerate}[label=(\alph*),topsep=2pt,itemsep=1pt,partopsep=0pt,parsep=0pt]
    \item \relax 45
    \item \relax 51
    \item \relax 63
    \item \relax 57
\end{enumerate}
\vspace{2mm}

\noindent\textbf{Q46.} \hfill {\small\color{LPUOrange}\textbf{[while loop | Hard]}}\\[2pt]
How many times does the loop body execute? 

\begin{pycode}
n = 100\\
count = 0\\
while n > 1:\\
n = n // 2\\
count += 1\\
print(count)
\end{pycode}

\begin{enumerate}[label=(\alph*),topsep=2pt,itemsep=1pt,partopsep=0pt,parsep=0pt]
    \item \relax 7
    \item \relax 6
    \item \relax 5
    \item \relax 8
\end{enumerate}
\vspace{2mm}

\noindent\textbf{Q47.} \hfill {\small\color{LPUOrange}\textbf{[nested loops | Hard]}}\\[2pt]
What is printed? 

\begin{pycode}
count = 0\\
for i in range(4):\\
for j in range(i):\\
count += 1\\
print(count)
\end{pycode}

\begin{enumerate}[label=(\alph*),topsep=2pt,itemsep=1pt,partopsep=0pt,parsep=0pt]
    \item \relax 16
    \item \relax 12
    \item \relax 10
    \item \relax 6
\end{enumerate}
\vspace{2mm}

\noindent\textbf{Q48.} \hfill {\small\color{LPUOrange}\textbf{[break and continue | Hard]}}\\[2pt]
What is the output? 

\begin{pycode}
result = []\\
for i in range(1, 10):\\
if i \% 2 == 0:\\
continue\\
if i > 6:\\
break\\
result.append(i)\\
print(result)
\end{pycode}

\begin{enumerate}[label=(\alph*),topsep=2pt,itemsep=1pt,partopsep=0pt,parsep=0pt]
    \item \relax [1, 2, 3, 4, 5]
    \item \relax [1, 3, 5]
    \item \relax [1, 3, 5, 7]
    \item \relax [1, 3, 5, 7, 9]
\end{enumerate}
\vspace{2mm}

\noindent\textbf{Q49.} \hfill {\small\color{LPUOrange}\textbf{[break and continue | Hard]}}\\[2pt]
What does this print? 

\begin{pycode}
for i in range(3):\\
for j in range(3):\\
if j == 1:\\
break\\
print(i, j)
\end{pycode}

\begin{enumerate}[label=(\alph*),topsep=2pt,itemsep=1pt,partopsep=0pt,parsep=0pt]
    \item \relax Prints \texttt{0 0}, \texttt{1 0}, \texttt{2 0}
    \item \relax Prints \texttt{0 0} only
    \item \relax Prints \texttt{0 0}, \texttt{0 1}, \texttt{0 2}
    \item \relax Prints all nine pairs
\end{enumerate}
\vspace{2mm}

\noindent\textbf{Q50.} \hfill {\small\color{LPUOrange}\textbf{[random numbers | Hard]}}\\[2pt]
Given \texttt{import random; random.seed(0)}, which statement is true about \texttt{random.randint(1, 6)}?

\begin{enumerate}[label=(\alph*),topsep=2pt,itemsep=1pt,partopsep=0pt,parsep=0pt]
    \item \relax It excludes the upper bound 6
    \item \relax It returns values from 1 to 5 only
    \item \relax It returns floats between 1 and 6
    \item \relax It can return any integer from 1 to 6 inclusive
\end{enumerate}
\vspace{2mm}

\noindent\textbf{Q51.} \hfill {\small\color{LPUOrange}\textbf{[random numbers | Hard]}}\\[2pt]
What is the possible range of values produced by \texttt{random.randrange(2, 20, 4)}?

\begin{enumerate}[label=(\alph*),topsep=2pt,itemsep=1pt,partopsep=0pt,parsep=0pt]
    \item \relax {2, 6, 10, 14, 18}
    \item \relax {4, 8, 12, 16}
    \item \relax {2, 6, 10, 14, 18, 20}
    \item \relax Any integer from 2 to 20
\end{enumerate}
\vspace{2mm}

\noindent\textbf{Q52.} \hfill {\small\color{LPUOrange}\textbf{[function calls | Hard]}}\\[2pt]
What is the output? 

\begin{pycode}
def g(x):\\
return x * 2\\
~\\
def h(x):\\
return x + 3\\
~\\
print(g(h(g(2))))
\end{pycode}

\begin{enumerate}[label=(\alph*),topsep=2pt,itemsep=1pt,partopsep=0pt,parsep=0pt]
    \item \relax 11
    \item \relax 14
    \item \relax 20
    \item \relax 10
\end{enumerate}
\vspace{2mm}

\noindent\textbf{Q53.} \hfill {\small\color{LPUOrange}\textbf{[type conversion and coercion | Hard]}}\\[2pt]
What is the result of the following? 

\begin{pycode}
print(int("10", 2) + int("10", 8) + int("10", 16))
\end{pycode}

\begin{enumerate}[label=(\alph*),topsep=2pt,itemsep=1pt,partopsep=0pt,parsep=0pt]
    \item \relax 30
    \item \relax 10
    \item \relax 26
    \item \relax 42
\end{enumerate}
\vspace{2mm}

\noindent\textbf{Q54.} \hfill {\small\color{LPUOrange}\textbf{[type conversion and coercion | Hard]}}\\[2pt]
What is the output? 

\begin{pycode}
print(int(True) + float("3.5") + int("7"))
\end{pycode}

\begin{enumerate}[label=(\alph*),topsep=2pt,itemsep=1pt,partopsep=0pt,parsep=0pt]
    \item \relax 10.5
    \item \relax 10
    \item \relax 11.5
    \item \relax 11
\end{enumerate}
\vspace{2mm}

\noindent\textbf{Q55.} \hfill {\small\color{LPUOrange}\textbf{[math functions | Hard]}}\\[2pt]
What does the following evaluate to? 

\begin{pycode}
import math\\
print(math.floor(math.sqrt(50)) + math.ceil(math.log2(10)))
\end{pycode}

\begin{enumerate}[label=(\alph*),topsep=2pt,itemsep=1pt,partopsep=0pt,parsep=0pt]
    \item \relax 10
    \item \relax 11
    \item \relax 12
    \item \relax 13
\end{enumerate}
\vspace{2mm}

\noindent\textbf{Q56.} \hfill {\small\color{LPUOrange}\textbf{[math functions | Hard]}}\\[2pt]
What is the output? 

\begin{pycode}
import math\\
print(math.gcd(48, 36), math.factorial(4) \% 7)
\end{pycode}

\begin{enumerate}[label=(\alph*),topsep=2pt,itemsep=1pt,partopsep=0pt,parsep=0pt]
    \item \relax 12 3
    \item \relax 6 4
    \item \relax 12 4
    \item \relax 6 3
\end{enumerate}
\vspace{2mm}

\noindent\textbf{Q57.} \hfill {\small\color{LPUOrange}\textbf{[adding new function | Hard]}}\\[2pt]
What is printed? 

\begin{pycode}
def make\_counter():\\
count = 0\\
def inc():\\
nonlocal count\\
count += 1\\
return count\\
return inc\\
~\\
c = make\_counter()\\
print(c(), c(), c())
\end{pycode}

\begin{enumerate}[label=(\alph*),topsep=2pt,itemsep=1pt,partopsep=0pt,parsep=0pt]
    \item \relax 1 1 1
    \item \relax 1 2 3
    \item \relax 0 1 2
    \item \relax 3 3 3
\end{enumerate}
\vspace{2mm}

\noindent\textbf{Q58.} \hfill {\small\color{LPUOrange}\textbf{[parameters and argument | Hard]}}\\[2pt]
What is the output? 

\begin{pycode}
def f(a, b=[]):\\
b.append(a)\\
return b\\
~\\
print(f(1))\\
print(f(2))\\
print(f(3, []))
\end{pycode}

\begin{enumerate}[label=(\alph*),topsep=2pt,itemsep=1pt,partopsep=0pt,parsep=0pt]
    \item \relax [1] then [1, 2] then [1, 2, 3]
    \item \relax [1] then [2] then [3]
    \item \relax [1, 2, 3] repeated three times
    \item \relax [1] then [1, 2] then [3]
\end{enumerate}
\vspace{2mm}

\noindent\textbf{Q59.} \hfill {\small\color{LPUOrange}\textbf{[recursion and its use | Hard]}}\\[2pt]
What does \texttt{mystery(4)} return? 

\begin{pycode}
def mystery(n):\\
if n <= 1:\\
return 1\\
return n * mystery(n - 1)
\end{pycode}

\begin{enumerate}[label=(\alph*),topsep=2pt,itemsep=1pt,partopsep=0pt,parsep=0pt]
    \item \relax 10
    \item \relax 12
    \item \relax 24
    \item \relax 4
\end{enumerate}
\vspace{2mm}

\noindent\textbf{Q60.} \hfill {\small\color{LPUOrange}\textbf{[recursion and its use | Hard]}}\\[2pt]
How many times is \texttt{fib} called (including the initial call) when computing \texttt{fib(4)}? 

\begin{pycode}
def fib(n):\\
if n < 2:\\
return n\\
return fib(n - 1) + fib(n - 2)
\end{pycode}

\begin{enumerate}[label=(\alph*),topsep=2pt,itemsep=1pt,partopsep=0pt,parsep=0pt]
    \item \relax 5
    \item \relax 9
    \item \relax 7
    \item \relax 11
\end{enumerate}
\vspace{2mm}


\newpage
\section{Section B: Comprehensive Subjective Problem with Detailed Analytical Rubric}

\begin{subjectiveq}{2}{1}
Design a complete recursive and iterative implementation of the binary search algorithm and memoized Fibonacci generator with call-stack analysis:
\begin{enumerate}[label=(\alph*)]
    \item Implement \texttt{binary\_search(arr, low, high, target)} recursively.
    \item Provide the formal recurrence relation for execution time and solve via Master Theorem.
    \item Implement \texttt{fibonacci\_memo(n, cache=None)} to evaluate $F(n)$ in $\mathcal{O}(n)$ time and compare stack-depth against naive recursion $\mathcal{O}(2^n)$.
\end{enumerate}
\end{subjectiveq}

\begin{solbox}
\textbf{Recursive Implementation \& Stack Analysis:}
\begin{pycode}
def binary\_search(arr, low, high, target):\\
\quad if low > high:\\
\quad \quad return -1\\
\quad mid = low + (high - low) // 2\\
\quad if arr[mid] == target:\\
\quad \quad return mid\\
\quad elif arr[mid] > target:\\
\quad \quad return binary\_search(arr, low, mid - 1, target)\\
\quad else:\\
\quad \quad return binary\_search(arr, mid + 1, high, target)\\
~\\
def fibonacci\_memo(n, cache=None):\\
\quad if cache is None:\\
\quad \quad cache = \{0: 0, 1: 1\}\\
\quad if n not in cache:\\
\quad \quad cache[n] = fibonacci\_memo(n - 1, cache) + fibonacci\_memo(n - 2, cache)\\
\quad return cache[n]
\end{pycode}
\textbf{Recurrence Relation:} For binary search, $T(n) = T(n/2) + \mathcal{O}(1)$. By Master Theorem ($a=1, b=2, d=0$), since $\log_2(1) = 0 = d$, $T(n) = \Theta(\log n)$. For memoized Fibonacci, each state $0 \dots n$ is visited exactly once, reducing exponential calls $2^n$ to linear time $\mathcal{O}(n)$ with auxiliary call stack depth $n$.
\end{solbox}

\begin{examinertip}
\begin{itemize}[leftmargin=1.5em, itemsep=2pt]
    \item \textbf{Algorithmic Formulation (3 Marks):} Correct problem decomposition, function signatures, and parameter documentation.
    \item \textbf{Syntactic Rigor \& Implementation (4 Marks):} Valid Python syntax, proper error handling, defensive input validation, and avoidance of anti-patterns.
    \item \textbf{Complexity \& Edge Analysis (3 Marks):} Accurate time and space asymptotic analysis $\mathcal{O}(n)$ and testing edge conditions (null inputs, zero division, empty collections).
\end{itemize}
\end{examinertip}
